% 330_Operationen_T4_Abgrenzung_En_ch.tex
% Doc. 330, FFGFT corpus, August 2026 — chapter file
% Included by Sources/wr_standalone_A4/330_Operationen_T4_Abgrenzung_En.tex

\providecommand{\K}{}\renewcommand{\K}{\textbf{[K]}}
\providecommand{\B}{}\renewcommand{\B}{\textbf{[B]}}
\renewcommand{\S}{\textbf{[S]}}
\providecommand{\SETZ}{}\renewcommand{\SETZ}{\textbf{[POSTULATE]}}
\providecommand{\E}{}\renewcommand{\E}{\textbf{[E]}}
\providecommand{\lP}{}\renewcommand{\lP}{\ell_{\mathrm{P}}}
\providecommand{\Tmass}{}\renewcommand{\Tmass}{\tilde{T}}

\definecolor{ffgftblue}{RGB}{30,90,160}
\definecolor{warncolor}{RGB}{140,0,40}
\tcbset{
  base/.style={breakable,enhanced,boxrule=0.6pt,arc=3pt,
    fonttitle=\bfseries\small,coltitle=white,top=4pt,bottom=4pt,left=6pt,right=6pt},
  ffgftbox/.style={base,colback=blue!5,colframe=ffgftblue,colbacktitle=ffgftblue},
  warnbox/.style={base,colback=red!5,colframe=warncolor,colbacktitle=warncolor}
}

\section*{Abstract}

This document consolidates the demarcation of FFGFT against discrete
emergence models (here: HLV by Krüger, August 2026) and clarifies
which operations mediate between $T^4$ and the observable space
$\mathbb{R}^3\times\mathbb{R}_t$.

Central claim: The failure of discrete lattice models (class-X4
overlaps, missing operator specificity, open continuum limit) is not
a surprise, but a structural consequence of treating geometric
auxiliary constructs as ontological building blocks. FFGFT avoids
these failures not by choosing a different lattice, but because the
fundamental level $T^4$ is compact and geometrically complete —
and the three operations leading from it to the observable space
have strictly different logical status.

The three operations (Doc.~270 \K):
\begin{enumerate}
  \item \textbf{Type~I} — Duality decompactification:
    $S^1_m \to \mathbb{R}_t$ via $\Tmass\cdot m = 1$.
    Information-preserving (embedding, not projection).
  \item \textbf{Type~II} — Geometric projection:
    $D3 \to D2 \to D1$. Lossy, irreversible.
    The step $D3 \to D2$ is identified as holographic projection.
  \item \textbf{Type~III} — Representational translation:
    $T^4 \leftrightarrow \mathcal{H}$ (Hilbert space/matrix).
    Bijective, lossless — upward movement, no geometric
    dimension reduction.
\end{enumerate}

% ============================================================
\section{Starting point: the HLV audit as case study}
\label{sec:hlv}

Marcel Krüger subjected his own Helix-Light-Vortex model (HLV) to
four prospectively frozen mathematical tests in \emph{Local-to-Global
Compatibility on Aperiodic Cut-and-Project Carriers} (v1.4,
August 2026). All key claims fail or remain open.

\begin{enumerate}
  \item \textbf{Local-to-global compatibility:}
    5,455 genuine volume overlaps (class~X4) among 3,384 projected
    tetrahedra. Six witnesses certified with 120-digit arithmetic.
    \emph{Failure.} \E
  \item \textbf{Operator specificity:}
    An ordinary periodic cubic lattice (CUBIC) outperforms HLV in
    spectral multiplet score (GAE-001A). \emph{Rejected.} \E
  \item \textbf{Geometry-to-constitutive law:}
    Pure carrier geometry does not determine physical densities,
    masses, or couplings (Proposition~2.2). \emph{Trivially negative.} \E
  \item \textbf{Finite-to-continuum:}
    A bounded graph has bounded spectrum, no UV limit.
    Mosco/resolvent convergence unproven. \emph{Open.} \E
\end{enumerate}

\begin{tcolorbox}[warnbox,title={Corollary (Krüger 4.5) \E}]
  Local validity of an operator does not imply global embeddability
  of the underlying carrier.
\end{tcolorbox}

\subsection{The structural cause}

HLV attempts to use geometric bodies (3D tetrahedra from
$\mathbb{R}^6$ projection) as ontological building blocks of
physical space. It thereby inherits all four problems.
The remedy is not a \emph{better} lattice — it is a fundamental
carrier that is geometrically complete and closed from the outset.

% ============================================================
\section{The FFGFT fundamental level: $T^4$ with D4 lattice}
\label{sec:t4}

\subsection{Structure of $T^4$}

The fundamental level of FFGFT is the compact 4-torus
\[
  T^4 = S^1_x \times S^1_y \times S^1_z \times S^1_m
\]
with three spatial circles $S^1_{x,y,z}$ (radius $R_H$,
quasi-continuous) and one mass circle $S^1_m$
(radius $R_m \sim \xi\cdot\lP$, discrete).
The fourth direction $S^1_m$ is not a mass dimension alongside
a separate time, but the substitute-time direction itself.
Mass and time are two readings of \emph{the same} dimension;
$\Tmass\cdot m = 1$ is the translation dictionary (Doc.~270 \K,
Doc.~314 \K). \SETZ

The lattice on $T^4$ is the D4 lattice (Doc.~314 \K):
kissing number~24, $|\text{Aut}(D_4)| = 1152$, triality as
built-in $\mathbb{Z}_3$ symmetry (16 order-3 elements generate
the $T^4/\mathbb{Z}_3$ orbifold with 9 fixed points).
Negative result with theorem character: $5\nmid 1152$, so
$\text{Aut}(D_4)$ contains no order-5 element \K.

\subsection{Why $T^4$ is flat — intrinsic vs.\ extrinsic curvature \K}
\label{sec:flat}

A torus carries \emph{extrinsic} curvature when embedded in a
higher-dimensional ambient space. The \emph{intrinsic} Gaussian
curvature, however, sums to zero pointwise — positive curvature
(outer side) and negative curvature (inner side) cancel.
By the Theorema Egregium, intrinsic curvature is an isometric
invariant, independent of embedding.

\begin{tcolorbox}[ffgftbox,title={Intrinsically flat torus \K}]
  $T^4$ is intrinsically flat: an interior observer measures no
  curvature. Angle sums are $\pi$, parallel geodesics remain
  parallel. $T^4$ is locally isometric to $\mathbb{R}^4$ —
  with global boundary identification, but without local
  curvature corrections.
\end{tcolorbox}

Two direct consequences for FFGFT:

\textbf{(a) D4 lattice through topology, not curvature.}
What the D4 lattice carries into the Hilbert space is
topological data: winding numbers, mode content, triality.
The distinction $\mathbb{Z}^4$/$D_4$ is spectral, not geometric
(Doc.~314 \K: ``The torus is flat anyway — curvature was never
the difference'').

\textbf{(b) Fourier analysis without curvature corrections.}
On the flat torus, Fourier modes $e^{ik\cdot x}$ form the natural
basis; Laplacian eigenvalues are $|k|^2$ without curvature
correction. This is the basis for the bijective, lossless
Type~III translation $T^4 \leftrightarrow L^2(T^4)\otimes\mathbb{C}^3$:
Fourier decomposition on a flat torus is a complete isometric map.

This structurally distinguishes FFGFT from approaches like HLV,
which embed geometric bodies into $\mathbb{R}^3$ and thereby
generate overlap problems: the embedding question does not arise
in FFGFT, because $T^4$ as an intrinsically flat, self-contained
carrier requires no embedding into an outer space.

\subsection{How information is encoded in the D4 lattice \K}
\label{sec:kodierung}

The geometry does not disappear when moving to the Hilbert space
— it migrates into spectrum and fibre (Doc.~314 \K):

\textbf{Spectrum = theta series.}
The degeneracies of the Laplacian on flat $T^4$ are exactly the
theta series of the D4 lattice: $\theta_{D_4}(q) = 1+24q^2+24q^4+\ldots$
The kissing number 24 becomes the degeneracy of the first excitation
$|k|^2=2$. Two flat tori ($\mathbb{Z}^4$ vs.\ $D_4$) are
distinguishable by spectrum alone, without curvature \K.

\textbf{Fibre = triality.}
Triality is built in, not appended \K: $|\text{Aut}(D_4)| = 1152$,
quotient to $W(D_4)$ is $6 = |S_3|$. The 16 order-3 elements
generate the $\mathbb{Z}_3$ orbifold and act on every mode orbit
as a $\mathbb{C}^3$ circulant. The fibre $\mathbb{C}^3$ follows
from the lattice automorphism structure, not an additional postulate.

\textbf{Winding numbers = topology.}
The reciprocity relation $\lambda_4\cdot m = 2\pi$ holds mode-by-mode
exactly \K: winding number $=$ mass number.
Winding numbers are topological — invariant under deformation;
fractal corrections accumulate only when unrolled, not when rolled up.

In one sentence (Doc.~314): ``In the Hilbert space the lattice
becomes flat, but not without consequence — what geometry carried,
spectrum and fibre carry, and the division of labour between them
is exactly determinable.''

\subsection{Where $\Tmass\cdot m = 1$ appears in the Hilbert space \K}
\label{sec:duality-hilbert}

The time-mass duality appears in $\mathcal{H} = L^2(T^4)\otimes\mathbb{C}^3$
not as time evolution, but in three complementary forms:

\textbf{(a) As the fourth winding number $n_4$.}
The multi-index $n=(n_1,n_2,n_3,n_4)\in\mathbb{Z}^4$ counts windings
in all four torus directions \B\ (Doc.~322).
The mass dimension $S^1_m$ carries the winding number $n_4$.
The reciprocity relation $\lambda_4\cdot m = 2\pi$ holds mode-by-mode
exactly \K\ (Doc.~314): winding number $=$ mass number.
$\Tmass\cdot m = 1$ is thus the Kaluza-Klein relation of the fourth
torus direction, readable as $n_4$ quantum number in the spectrum.

\textbf{(b) As the mass operator eigenvalue equation.}
\begin{equation}
  \hat{M}\,|n_\theta,n_\phi,p\rangle = m_i\,|n_\theta,n_\phi,p\rangle
  \quad\textbf{[K]}
\end{equation}
The eigenvalues are the particle masses (Doc.~322, Doc.~320).
The duality appears as the spectrum of the operator.

\textbf{(c) As radius-class separation in the spectrum.}
The splitting at $R_4\neq R_{123}$ makes $\Tmass\cdot m = 1$
directly readable \K\ (Doc.~314): the mass dimension generates
a large mode spacing $\sim 1/R_m$, the spatial directions a
quasi-continuous spacing $\sim 1/R_H$.

\begin{tcolorbox}[ffgftbox,title={What does \emph{not} appear in the Hilbert space}]
  Time $\mathbb{R}_t$ itself. It is the result of the Type-I
  decompactification $S^1_m\to\mathbb{R}_t$ — a geometric downward
  operation \emph{before} the Type-III translation.
  In the Hilbert space $L^2(T^4)$, only the torus lives;
  time appears only when leaving the Hilbert space and embedding
  the spectrum into $\mathbb{R}_t$.
  $\Tmass\cdot m=1$ resides in the Hilbert space as $n_4$ winding
  number and mass operator eigenvalue structure, not as time evolution.
\end{tcolorbox}

\subsection{Where the fractal property resides in the Hilbert space \K\ / \B}
\label{sec:fractal-hilbert}

The fractal property of FFGFT appears in the Hilbert space at
three precisely separable locations:

\textbf{(a) In the operator $\hat{F}$ itself — as self-similarity structure \B.}
\begin{equation}
  \hat{F} = \bigoplus_{n=1}^{N} S_n\circ L_n,\quad
  S_n = r_n U_n,\quad r_n = \xi^n\in(0,1)
\end{equation}
The scale operators $S_n$ implement discrete scale transformations;
each level $n$ is a contraction of the previous by $r_n$ —
this is self-similarity as an operator.
$\hat{F}$ encodes the sub-Planck geometry as a hierarchical
convolution of scale levels down to the minimum length scale
$L_0 = \xi\cdot\lP$ (Doc.~180, 327 \K).

\textbf{(b) In the fractal dimension $D_f = 3-\xi$ — as spectral property \K.}
\begin{equation}
  D_f = \inf\!\left\{s:\sum_i\lambda_i^s<\infty\right\}
  = 3-\xi \approx 2.9999
  \quad\textbf{[K]}
\end{equation}
$D_f$ is defined via the Hausdorff measure of the spectrum of
$\hat{F}$ (Doc.~322): the spectrum on $L^2(T^4)$ has a fractal
fine structure. The value is spatial ($3-\xi$, three dimensions),
not four-dimensional.
Local action: $D_f^\text{space} = 3-\xi$ with $6.7\times10^{-5}$
relative (Doc.~314 \K).

\textbf{(c) In the interface factor $K_\text{frac}$ — as scale bridge \K.}
\begin{equation}
  K_\text{frac} = 1-100\xi = \tfrac{74}{75} \approx 0.9867
\end{equation}
$K_\text{frac}$ is cumulative (Doc.~133): it accumulates over the
RG running across 17 orders of magnitude and acts at the transition
from geometric (bare) to measured (SI) mass:
$m_i^\text{SI} = K_\text{frac}\cdot C_\text{conv}\cdot m_i^\text{bare}$.

\begin{tcolorbox}[ffgftbox,title={Decisive separation \K}]
  In the \emph{rolled-up} $T^4$ (Hilbert space) there are
  \textbf{no} fractal corrections —
  winding numbers are purely topological (Doc.~314 \K).
  The fractal property appears only when unrolling
  (Type~I: $S^1_m\to\mathbb{R}_t$) as an accumulated defect.
  In the Hilbert space itself, $\hat{F}$ is a bounded
  self-adjoint operator with real spectrum \B\ —
  the fractality resides in the scale hierarchy of the operator
  structure ($r_n = \xi^n$), not in the carrier $L^2(T^4)$.
\end{tcolorbox}

% ============================================================
\section{No embedding problem in $T^4$}
\label{sec:noembed}

$T^4$ is compact, complete, and geometrically self-contained.
There is no embedding problem because $T^4$ does not need to be
\emph{embedded} in an outer space: it \emph{is} the fundamental
level. The class-X4 errors of HLV arise because projection elements
must be embedded into an external $\mathbb{R}^3$. In FFGFT
this question does not arise.

\subsection{Type~III: representational translation upward}

The translation $T^4\leftrightarrow\mathcal{H} = L^2(T^4)\otimes\mathbb{C}^3$
(Doc.~230, 314, 322) is a bijective, lossless upward movement
(Type~III, Doc.~270 \K).
It adds no geometric extra dimensions, but expresses the same
$T^4$ content in a higher-dimensional language.
The spectrum of $\hat{F}$ (Doc.~322, self-adjoint \B, Doc.~327)
has the particle masses as eigenvalues;
$\xi = \lambda_{\min}(\hat{F}_{D_4})$ is the smallest eigenvalue \B.

% ============================================================
\section{The three operations from $T^4$ to the observable space}
\label{sec:operations}

\subsection{Type~I: duality decompactification}
\[
  S^1_m\;\xrightarrow{\;\Tmass\cdot m=1\;}\;\mathbb{R}_t
\]
The unrolling of the compact mass dimension into the non-compact
time direction is an \emph{embedding}, not a projection.
No mode content is lost; time arises as an additional non-compact
direction \K.

\subsection{Type~II: geometric projection (lossy)}

The further steps $D3\to D2\to D1$ are genuine geometric
projections of spatial directions: lossy, irreversible.
The step $D3\to D2$ is identified as holographic projection
(Doc.~270 \K). It is not fundamental.

\begin{tcolorbox}[ffgftbox,title={Classification of the holographic step}]
  The step $D3\to D2$ is \emph{not} the fundamental level of FFGFT.
  It is a lossy observer projection — physically relevant for
  surface phenomena (black hole membrane, Doc.~325), but not
  for the derivation of masses and couplings anchored on $T^4$.
  The two-dimensional projection surface is not a fundamental object.
\end{tcolorbox}

\subsection{Three-way comparison}

\begin{table}[h!]\small
\begin{tabular}{@{}llll@{}}
\toprule
Type & Operation & Direction & Status \\
\midrule
I   & Duality decompactification $S^1_m\to\mathbb{R}_t$
    & geometrically down & information-preserving \K \\
II  & Geometric projection $D3\to D2\to D1$
    & geometrically down & lossy \K \\
III & Representational translation $T^4\leftrightarrow\mathcal{H}$
    & representationally up & bijective, lossless \K \\
\bottomrule
\end{tabular}
\end{table}

% ============================================================
\section{Open questions in the flat Hilbert space — status}
\label{sec:hilbert}

\subsection{Question 1: Haar measure and D4 lattice \B}

The canonical Haar measure on $T^4$ is compatible with the D4
spectral structure (scripts 330-1, all 7/7):
Fourier modes are exactly orthogonal; D4 contains only even norm
shells; $\text{Aut}(D_4)\subset\text{Iso}(T^4)$ leaves the Haar
measure invariant. The $\mathbb{Z}_3$ orbifold breaks translation
invariance and splits the 24-fold degeneracy to $9\oplus8\oplus4\oplus2\oplus1$
(Doc.~314 \K) — but only as fine structure $\sim\xi\approx10^{-4}$,
below any measurement threshold. \B

\subsection{Question 2: Orbifold fixed points \K\ / \S}

Nine fixed points of $T^4/\mathbb{Z}_3$ follow from the Lefschetz
fixed-point formula (script 330-2, B1):
$|\det(1-A)| = |(1-\omega)(1-\omega^2)|^2 = 9$.
Neutrino localisation: $\nu_1\leftrightarrow F_2$,
$\nu_2\leftrightarrow F_5$, $\nu_3\leftrightarrow F_7$ \K\ (Doc.~322).
What remains \S: back-reaction of Hawking emission on the fixed-point
structure during $M\to M_\text{coll}=\mP/\sqrt{2}$ (explicitly open,
Doc.~325).

\subsection{Question 3: Fractal weighting $r_n$ and $L^2(T^4)$ \B}

$r_n = \xi^n$ is position-independent: $\hat{F}$ does \emph{not}
break the translation invariance of $L^2(T^4)$ — only the
$\mathbb{Z}_3$ orbifold does (and only as fine structure $\sim\xi$).
$\hat{F}$ is bounded ($\|\hat{F}\|\leq\sum r_n<\infty$) and
self-adjoint on $\mathcal{H}$ \B\ (Doc.~327). \B

\begin{tcolorbox}[ffgftbox,title={Status of the three questions}]
  Question 1 (Haar/D4): \B\ — fully resolved.\\
  Question 2 (fixed points): \K\ for 9 fixed points and
  $\nu$-localisation; \S\ for back-reaction during evaporation
  (explicitly open, Doc.~325).\\
  Question 3 ($r_n$/$\hat{F}$): \B\ — $\hat{F}$ bounded and
  self-adjoint; $r_n$ does not break translation invariance.
\end{tcolorbox}

% ============================================================
\section{Why Krüger's four problems do not arise in FFGFT}
\label{sec:comparison}

\begin{enumerate}
  \item \textbf{Class-X4 overlaps:} arise because HLV embeds projection
    elements into $\mathbb{R}^3$. In FFGFT, $T^4$ is the fundamental
    level; no embedding into an outer space occurs.
  \item \textbf{Missing operator specificity:} In FFGFT the spectrum
    is determined by the D4 lattice with built-in triality.
    The negative result $5\nmid|\text{Aut}(D_4)|$ is itself a specific
    theorem about the spectrum \K.
  \item \textbf{Geometry without physics:} $T^4$ alone also does not
    fix coupling constants. The physics comes from the combination
    $T^4$-spectrum $+$ $\xi = \lambda_{\min}(\hat{F}_{D_4})$.
  \item \textbf{Continuum limit:} No limit problem, because Type~I
    is an embedding (no limit passage needed) and Type~III is
    bijective (no information loss). Only Type~II is lossy —
    but this is known and documented (Doc.~270).
\end{enumerate}

% ============================================================
\section{What remains from the comparative discussion}
\label{sec:balance}

\textbf{Structurally secured:} The four HLV problems do not arise
in FFGFT — due to $T^4$ as intrinsically flat, self-contained
carrier, bijective Type-III translation, and Type-I embedding
$S^1_m\hookrightarrow\mathbb{R}_t$.

\textbf{Correct from the comparative discussion:}
The holographic projection $D3\to D2$ (Type~II) is lossy and not
fundamental; mass and time are not encoded in the 2D surface but
arise through Type~I ($\Tmass\cdot m=1$) and the recursion $\xi^n$.

\textbf{Correction vs.\ uncritical adoption:}
There is no ``two-dimensional projection surface as fundamental
object'' in FFGFT. The corpus has a clear ontological hierarchy:
$T^4$ as primary reality (Doc.~152 \SETZ), not a 2D surface.
$\Tmass\cdot m=1$ is the Kaluza-Klein relation of the mass dimension,
not a ``reading principle on a screen''.

\vspace{4pt}\noindent
\textbf{Verification scripts:}
\texttt{pruef\_330\_1\_haar\_d4.py} (7/7),
\texttt{pruef\_330\_2\_fixpunkte\_rn.py} (9/9).

% ============================================================
\section*{Status overview}

\begin{table}[h!]\small
\begin{tabular}{@{}lc@{}}
\toprule
Element & Status \\
\midrule
$T^4 = S^1_x\times S^1_y\times S^1_z\times S^1_m$ as fundamental level & \SETZ \\
D4 lattice on $T^4$; kissing number~24; $|\text{Aut}(D_4)|=1152$ & \K \\
$T^4$ intrinsically flat; spectral structure from topology & \K \\
Information in spectrum (theta series) and fibre (triality) & \K \\
$\Tmass\cdot m=1$ as $n_4$ winding number and mass operator eigenvalue & \K \\
Fractal property in $\hat{F}$ structure ($r_n=\xi^n$), not in $L^2(T^4)$ & \B \\
$D_f = 3-\xi$; no fractal corrections in rolled-up $T^4$ & \K \\
Three operation types (Doc.~270): Type~I/II/III & \K \\
Type~III bijective ($T^4\leftrightarrow\mathcal{H}$, Doc.~322/314) & \K \\
HLV class-X4 errors structurally excluded (no embedding problem) & \B \\
Continuum limit structurally resolved by Type~I (embedding) & \K \\
\bottomrule
\end{tabular}
\end{table}

\vspace{6pt}\noindent\small
\textcopyright\ 2026 Johann Pascher · \href{https://creativecommons.org/licenses/by/4.0/}{CC BY 4.0}
